% GENERATED FILE. Edit canonical lessons, then run npm run generate:books.
% canonical-source-hash: fnv1a64:05a2b01739587914
% canonical-lessons: GU-C147-hato, GU-C147-gaikale, GU-C147-gayo, GU-C147-thayo, GU-C147-modu, GU-W47-dda, GU-C147-past-practice

\chapter{What was and where I went}
\label{ch:gu-what-was-and-where-i-went}
\hlchaptermodality{147}

\begin{chapteropening}
\textbf{By the end of this chapter:} \emph{I can say where I was and where I went yesterday, with the past of to be, of going and coming, and of happening, and say that I came late.}

Tell a day in the past from morning to night, then tell your own yesterday in two sentences.
\end{chapteropening}

\section[\texorpdfstring{\gu{હતો} · \gu{હતી} · \gu{હતું} … (hato · hatī …)}{હતો · હતી · હતું … (hato · hatī …)}]{\gu{હતો} — \emph{was}, and an ending that listens to the subject}

\label{lesson:GU-C147-hato}



\begin{quote}
Ask \emph{do you live here?} Then say \emph{my brother} and \emph{my sister}: their endings \textbf{-\gu{ો}} and \textbf{-\gu{ી}} are about to come back. And the word for \emph{there}: \textbf{\gu{ત્યાં}}.
\end{quote}

\begin{grammarlens}[title={Grammar Lens: \gu{છે} becomes \gu{હતો}}]
\begin{quote}
\textbf{\gu{હું} \gu{ત્યાં} \gu{હતો}.} — \emph{hũ tyā̃ hato.} — I was there.
\end{quote}

In the present the copula changes for the person: \textbf{\gu{છું}}, \textbf{\gu{છે}}, \textbf{\gu{છો}}. In the past it stops caring about the person and starts caring about gender, with the endings \emph{my} uses:

\begin{quote}
\textbf{\gu{હું} \gu{ત્યાં} \gu{હતી}.} — \emph{hũ tyā̃ hatī.} — I was there, said by a woman.
\end{quote}
\end{grammarlens}

\subsection*{You'll want to know: \gu{હતું} and \gu{હતા}}
A neuter subject takes \textbf{\gu{હતું}}. More than one person, or respect, takes \textbf{\gu{હતા}}:

\begin{quote}
\textbf{\gu{ઘર} \gu{નાનું} \gu{હતું}.} — \emph{ghar nānũ hatũ.} — The house was small. \textbf{\gu{તેઓ} \gu{અહીં} \gu{હતા}.} — \emph{teo ahī̃ hatā.} — They were here.
\end{quote}

\begin{itemize}
\raggedright
  \item \textbf{\gu{હતો}} — he, or a man saying \emph{I};
  \item \textbf{\gu{હતી}} — she, or a woman saying \emph{I};
  \item \textbf{\gu{હતું}} — a neuter thing;
  \item \textbf{\gu{હતા}} — more than one, or respect.
\end{itemize}

\textbf{\gu{મારો} \gu{મારી} \gu{મારું} \gu{મારા}} and \textbf{\gu{હતો} \gu{હતી} \gu{હતું} \gu{હતા}} end alike, sign for sign.

\subsection*{Your turn}
\begin{itemize}
\raggedright
  \item \emph{Say it:} \emph{hũ tyā̃ hato}, or \emph{hatī}
  \item \emph{Say it:} \textbf{\gu{તે} \gu{અહીં} \gu{હતી}} — \emph{te ahī̃ hatī}, she was here
  \item \emph{Say it:} \textbf{\gu{અમે} \gu{ત્યાં} \gu{હતા}} — \emph{ame tyā̃ hatā}, we were there
  \item \emph{Choose:} the ending for \textbf{\gu{પાણી}}, a neuter noun: \textbf{\gu{પાણી} \gu{ગરમ} \gu{હતું}}
\end{itemize}

\begin{tcolorbox}[breakable,colback=teal!4,colframe=teal!35!black,title={Before you move on}]
What does the past copula agree with? (The subject's gender and number, not its person.) Which form goes with \textbf{\gu{ઘર}}? (\textbf{\gu{હતું}}.)
\end{tcolorbox}
\section[\texorpdfstring{\gu{ગઈકાલે} · \gu{સવારે} · \gu{ઘરે} (gaīkāle · savāre …)}{ગઈકાલે · સવારે · ઘરે (gaīkāle · savāre …)}]{\gu{ગઈકાલે}, \gu{સવારે}, \gu{ઘરે} — the -e of \emph{when} and \emph{where}}

\label{lesson:GU-C147-gaikale}



\begin{quote}
Say \emph{I was there}, then \emph{they were here}. And the word for \emph{yesterday}: \textbf{\gu{ગઈકાલ}}.
\end{quote}

\begin{grammarlens}[title={Grammar Lens: -e on a time}]
\textbf{\gu{ગઈકાલ}} names the day. To put something \emph{on} that day, Gujarati adds \textbf{-\gu{ે}}:

\begin{quote}
\textbf{\gu{ગઈકાલે} \gu{હું} \gu{ત્યાં} \gu{હતો}.} — \emph{gaīkāle hũ tyā̃ hato.} — Yesterday I was there.
\end{quote}

The same \textbf{-\gu{ે}} turns every part of the day, and every day, into a \emph{when}:

\begin{itemize}
\raggedright
  \item \textbf{\gu{આજ}} $\to$ \textbf{\gu{આજે}} — \emph{āje} — today;
  \item \textbf{\gu{આવતીકાલ}} $\to$ \textbf{\gu{આવતીકાલે}} — \emph{āvtīkāle} — tomorrow;
  \item \textbf{\gu{સવાર}} $\to$ \textbf{\gu{સવારે}} — \emph{savāre} — in the morning;
  \item \textbf{\gu{બપોર}} $\to$ \textbf{\gu{બપોરે}} — \emph{bapore} — at midday;
  \item \textbf{\gu{સાંજ}} $\to$ \textbf{\gu{સાંજે}} — \emph{sānje} — in the evening;
  \item \textbf{\gu{રવિવાર}} $\to$ \textbf{\gu{રવિવારે}} — \emph{ravivāre} — on Sunday.
\end{itemize}

\textbf{\gu{રાત}} is the one that changes shape: \textbf{\gu{રાત્રે}} (\emph{rātre}), \emph{at night}, keeps an older form with \textbf{\gu{ત્ર}}.
\end{grammarlens}

\subsection*{You'll want to know: -e on a place, and the two-way \gu{કાલે}}
\textbf{\gu{ઘર}} takes it too: \textbf{\gu{ઘરે}} (\emph{ghare}), \emph{at home}.

\textbf{\gu{કાલે}} has been in the book since \textbf{\gu{કાલે} \gu{મળીશું}}, \emph{see you tomorrow}. It also means \emph{yesterday}, and the verb decides which:

\begin{itemize}
\raggedright
  \item \textbf{\gu{કાલે} \gu{મળીશું}.} — a future verb, so \emph{tomorrow};
  \item \textbf{\gu{કાલે} \gu{હું} \gu{ઘરે} \gu{હતી}.} — a past verb, so \emph{yesterday}.
\end{itemize}

\subsection*{Your turn}
\begin{itemize}
\raggedright
  \item \emph{Say it:} \emph{gaīkāle hũ ghare hato}, or \emph{hatī}
  \item \emph{Say it:} \textbf{\gu{રાત્રે} \gu{તેઓ} \gu{ઘરે} \gu{હતા}} — \emph{rātre teo ghare hatā}
  \item \emph{Answer:} is \textbf{\gu{કાલે}} in \emph{kāle hũ ghare hatī} yesterday or tomorrow? (Yesterday: the verb is past.)
  \item \emph{Say it:} the answer with \textbf{\gu{નહીં}}: \textbf{\gu{આવતીકાલે} \gu{નહીં}, \gu{ગઈકાલે}} — \emph{āvtīkāle nahī̃, gaīkāle}
\end{itemize}

\begin{tcolorbox}[breakable,colback=teal!4,colframe=teal!35!black,title={Before you move on}]
What does \textbf{-\gu{ે}} do to \textbf{\gu{સવાર}}? (It makes the time a \emph{when}: \emph{in the morning}.) How do you know which day \textbf{\gu{કાલે}} means? (From the verb's tense.)
\end{tcolorbox}
\section[\texorpdfstring{\gu{ગયો} · \gu{ગઈ} · \gu{ગયું} … (gayo · gaī …)}{ગયો · ગઈ · ગયું … (gayo · gaī …)}]{\gu{હું} \gu{ગયો}, \gu{હું} \gu{ગઈ} — the past of moving}

\label{lesson:GU-C147-gayo}



\begin{quote}
Say \emph{yesterday I was at home}, and say which day \textbf{\gu{કાલે}} means in \textbf{\gu{કાલે} \gu{મળીશું}}. Then the two verbs of movement: \textbf{\gu{જવું}}, \emph{to go}, and \textbf{\gu{આવવું}}, \emph{to come}.
\end{quote}

\begin{grammarlens}[title={Grammar Lens: one word does the whole job}]
\begin{quote}
\textbf{\gu{ગઈકાલે} \gu{હું} \gu{બજાર} \gu{ગયો}.} — \emph{gaīkāle hũ bajār gayo.} — Yesterday I went to the market.
\end{quote}

\textbf{\gu{ગયો}} is a past of \textbf{\gu{જવું}} all on its own: no \textbf{\gu{છું}}, no \textbf{\gu{હતો}} beside it. It ends the way \textbf{\gu{હતો}} ends, and it agrees the same way:

\begin{itemize}
\raggedright
  \item \textbf{\gu{ગયો}} — a man; \textbf{\gu{ગઈ}} — a woman;
  \item \textbf{\gu{ગયું}} — something neuter; \textbf{\gu{ગયા}} — more than one, or respect.
\end{itemize}

\textbf{\gu{ગઈ}} is not new to you. It is the first half of \textbf{\gu{ગઈકાલ}}, \emph{yesterday}: the gone day. \textbf{\gu{ગયો}} goes back through Old Gujarati \emph{gayaü} to Sanskrit \emph{gata}, \textquotedblleft{}gone\textquotedblright{} (CDIAL 4008).
\end{grammarlens}

\subsection*{You'll want to know: \gu{આવ્યો}}
\textbf{\gu{આવવું}} builds its past on \textbf{\gu{આવ્ય}-}:

\begin{quote}
\textbf{\gu{તે} \gu{ગઈકાલે} \gu{ઘરે} \gu{આવ્યો}.} — \emph{te gaīkāle ghare āvyo.} — He came home yesterday. \textbf{\gu{તે} \gu{આવી}.} — \emph{te āvī.} — She came.
\end{quote}

The same four endings again: \textbf{\gu{આવ્યો}}, \textbf{\gu{આવી}}, \textbf{\gu{આવ્યું}}, \textbf{\gu{આવ્યા}}. The feminine sheds the \textbf{\gu{્ય}}.

\subsection*{Your turn}
Say these aloud:

\begin{itemize}
\raggedright
  \item \emph{hũ gayo}, \emph{hũ gaī}, \emph{ame gayā}
  \item \textbf{\gu{મારી} \gu{બહેન} \gu{સાંજે} \gu{આવી}} — \emph{mārī bahen sānje āvī}
  \item \textbf{\gu{તમે} \gu{બજાર} \gu{ગયા}, \gu{ખરું}?} — \emph{tame bajār gayā, kharũ?}, you went to the market, right?
\end{itemize}

\begin{tcolorbox}[breakable,colback=teal!4,colframe=teal!35!black,title={Before you move on}]
Which other words does the past of \emph{go} need? (None.) What does its ending agree with? (Whoever went.) Which old word holds \textbf{\gu{ગઈ}}? (\textbf{\gu{ગઈકાલ}}.)
\end{tcolorbox}
\section[\texorpdfstring{\gu{થયો} · \gu{થઈ} · \gu{થયું} … (thayo · thaī …)}{થયો · થઈ · થયું … (thayo · thaī …)}]{\gu{થયું} — \emph{it happened}}

\label{lesson:GU-C147-thayo}



\begin{quote}
Say \emph{I went}, then \emph{she came}. And the greeting from the second day: \textbf{\gu{તમને} \gu{મળીને} \gu{આનંદ} \gu{થયો}}, \emph{pleased to meet you}.
\end{quote}

\begin{grammarlens}[title={Grammar Lens: the past that was there from the start}]
\textbf{\gu{થયો}} is the past of \textbf{\gu{થવું}}, \emph{to happen, to become}. In \textbf{\gu{તમને} \gu{મળીને} \gu{આનંદ} \gu{થયો}} the joy is the subject (\emph{to you, meeting, joy happened}), and \textbf{\gu{આનંદ}} is masculine, so the verb says \textbf{\gu{થયો}}.

Change the subject and the ending follows it:

\begin{itemize}
\raggedright
  \item \textbf{\gu{આનંદ} \gu{થયો}} — joy happened: masculine;
  \item \textbf{\gu{મને} \gu{ચિંતા} \gu{થઈ}} — worry happened to me, \emph{I got worried}: feminine;
  \item \textbf{\gu{કામ} \gu{થયું}} — the work got done: neuter;
  \item \textbf{\gu{તેઓ} \gu{ખુશ} \gu{થયા}} — they became happy: more than one.
\end{itemize}

\textbf{\gu{થયો}} goes back to Sanskrit \emph{sthita}, \textquotedblleft{}standing\textquotedblright{}, which CDIAL lists as a past of \emph{to be} (CDIAL 13768).
\end{grammarlens}

\subsection*{Your turn}
Say these aloud:

\begin{itemize}
\raggedright
  \item \emph{mane ānand thayo}
  \item \textbf{\gu{મને} \gu{ચિંતા} \gu{થઈ}} — \emph{mane cintā thaī}
  \item \textbf{\gu{કામ} \gu{થયું}} — \emph{kām thayũ}, done
  \item \textbf{\gu{તેઓ} \gu{ખુશ} \gu{થયા}} — \emph{teo khuś thayā}
  \item \textbf{\gu{તેમને} \gu{આનંદ} \gu{થયો}} — \emph{temne ānand thayo}, they were pleased
  \item the reply to \emph{I got worried}: \textbf{\gu{ચિંતા} \gu{ન} \gu{કરો}} — \emph{cintā na karo}
\end{itemize}

\begin{tcolorbox}[breakable,colback=teal!4,colframe=teal!35!black,title={Before you move on}]
In \textbf{\gu{મને} \gu{ચિંતા} \gu{થઈ}}, what is the subject? (The worry; \emph{I} am \textbf{\gu{મને}}.) Why \textbf{\gu{થઈ}}? (\textbf{\gu{ચિંતા}} is feminine.)
\end{tcolorbox}
\section[\texorpdfstring{\gu{મોડું} (moḍũ)}{મોડું (moḍũ)}]{\gu{મોડું} (moḍũ) — late}

\label{lesson:GU-C147-modu}



\begin{quote}
Before the new one: say \emph{the work got done}, then \emph{he came}.
\end{quote}

\subsection*{You'll want to know: \gu{મોડું}}
\textbf{\gu{મોડું}} (\emph{moḍũ}) — \textquotedblleft{}late\textquotedblright{}.

It is the partner of \textbf{\gu{વહેલું}}, \emph{early}, and like the possessives it bends to agree:

\begin{quote}
\textbf{\gu{હું} \gu{મોડો} \gu{આવ્યો}.} — \emph{hũ moḍo āvyo.} — I came late. (a man) \textbf{\gu{તે} \gu{મોડી} \gu{આવી}.} — \emph{te moḍī āvī.} — She came late. \textbf{\gu{મોડું} \gu{થયું}.} — \emph{moḍũ thayũ.} — It got late.
\end{quote}

\begin{cousinweb}
\textbf{\gu{મોડું}} (\emph{moḍũ}) is an old word whose origin is not settled: CDIAL lists no Gujarati form for it, and no Sanskrit source has been shown. Its \textbf{-\gu{ું}} is the neuter ending of the adjective, which is why it turns into \textbf{\gu{મોડો}} and \textbf{\gu{મોડી}}.
\end{cousinweb}

\begin{grammarlens}[title={What you've built}]
One more word, and it arrives already in the past.
\end{grammarlens}

\subsection*{Your turn}
Say these aloud:

\begin{itemize}
\raggedright
  \item \emph{moḍũ}
  \item \emph{hũ moḍo āvyo}, or \emph{moḍī āvī}
  \item \emph{moḍũ thayũ}
  \item \textbf{\gu{તેનો} \gu{ભાઈ} \gu{મોડો} \gu{આવ્યો}} — \emph{teno bhāī moḍo āvyo}, her brother came late
  \item \textbf{\gu{હું} \gu{ક્યારેય} \gu{મોડો} \gu{નથી} \gu{આવતો}} — \emph{hũ kyārey moḍo nathī āvto}, I never come late
\end{itemize}

\begin{tcolorbox}[breakable,colback=teal!4,colframe=teal!35!black,title={Before you move on}]
What does \textbf{\gu{મોડું}} mean? (\emph{Late}.) Which letter in it have you read but never written? (\textbf{\gu{ડ}}, the next lesson.)
\end{tcolorbox}
\section[\texorpdfstring{\gu{ડ} (ḍa)}{ડ (ḍa)}]{\gu{ડ} — retroflex ḍa}

\label{lesson:GU-W47-dda}



\begin{quote}
\begin{itemize}
\raggedright
  \item \emph{Recall:} say \emph{late} — \textbf{R1}, one lesson back
  \item \emph{Recall:} say \textbf{\gu{ટ}}, \textbf{\gu{ઠ}} and \textbf{\gu{ઢ}}, the three tongue-curled letters you already write
\end{itemize}

\textbf{\gu{મોડું}} holds the one curled consonant your hand has not made yet.
\end{quote}

\subsection*{You'll want to know: the fourth corner}
Gujarati's retroflex row, with the tongue curled back, has a plain and a breathy letter for each voicing:

\noindent
\begin{tabularx}{\linewidth}{@{}>{\raggedright\arraybackslash}X>{\raggedright\arraybackslash}X>{\raggedright\arraybackslash}X@{}}
\toprule
\textbf{} & \textbf{plain} & \textbf{breathy} \\
\midrule
voiceless & \textbf{\gu{ટ}} \emph{ṭa} & \textbf{\gu{ઠ}} \emph{ṭha} \\
voiced & ? & \textbf{\gu{ઢ}} \emph{ḍha} \\
\bottomrule
\end{tabularx}

The missing corner is the plain voiced one, \emph{ḍa}, and \textbf{\gu{મોડું}} needs it.

\subsection*{Script — meet one form}
\begin{quote}
\textbf{\gu{ડ}} — \emph{ḍa}
\end{quote}

Keep \textbf{\gu{ડ}} visible. Begin at the upper right and sweep left across the high shoulder. Come down through the middle into the right side of the broad lower bowl, then carry the bowl round its base and finish at its lower-left end. One stroke, no lift.

\textbf{\gu{ઠ}} starts the same way, but its bowl curls back inward and closes. \textbf{\gu{ડ}} leaves its bowl open at the lower left. If the end turns back in, you have made a \textbf{\gu{ઠ}}.

\subsection*{Writing — notice and trace}
\hlblockfigure{\detokenize{figures/GU-W47-dda-filmstrip.pdf}}{How \gu{ડ} is written, stroke by stroke}

Finger-trace \textbf{\gu{ડ}} once, then air-trace it once. Notice the shoulder, the drop through the middle, and the open end of the bowl.

\subsection*{Writing — copy with the model}
Copy visible \textbf{\gu{ડ}} once. Compare only two things: is the bowl broad, and does it end open at the lower left? Repair one landmark if needed.

\subsection*{Writing — recall after hiding}
Hide the model and write retroflex \emph{ḍa} once. Uncover \textbf{\gu{ડ}}, compare, and stop after one repair.

\begin{tcolorbox}[breakable,colback=teal!4,colframe=teal!35!black,title={Before you move on}]
Read \textbf{\gu{ડ} \gu{ઠ} \gu{મોડું}} aloud. Which corner of the retroflex row is still missing? (\textbf{None}: with \textbf{\gu{ડ}} all four are in your hand.)

Sources: \href{https://www.t30apps.com/gujarati-alphabet-writing/}{Gujarati Alphabet Writing Practice} and the \href{https://www.unicode.org/charts/PDF/U0A80.pdf}{Unicode Gujarati chart}.
\end{tcolorbox}
\section[Practice]{Yesterday — a practice}

\label{lesson:GU-C147-past-practice}



\begin{quote}
Say \emph{I came late}.

\begin{itemize}
\raggedright
  \item \emph{Write it:} \textbf{\gu{ડ}} once from memory
\end{itemize}
\end{quote}

\begin{tcolorbox}[breakable,skin=enhanced,colback=orange!6,colframe=orange!45!black,boxrule=0.5pt,arc=1mm,left=6pt,right=6pt,top=4pt,bottom=4pt,fonttitle=\bfseries,title={Read this}]
\begin{quote}
\textbf{\gu{ગઈકાલે} \gu{સવારે} \gu{હું} \gu{ઘરે} \gu{હતી}.} \textbf{\gu{બપોરે} \gu{હું} \gu{બજાર} \gu{ગઈ}.} \textbf{\gu{સાંજે} \gu{મારો} \gu{ભાઈ} \gu{આવ્યો}.} \textbf{\gu{તે} \gu{મોડો} \gu{આવ્યો}.} \textbf{\gu{રાત્રે} \gu{અમે} \gu{ઘરે} \gu{હતા}. \gu{મને} \gu{આનંદ} \gu{થયો}.}
\end{quote}

Say these aloud:

\begin{itemize}
\raggedright
  \item the five lines aloud
  \item the same day as a man would tell it: \emph{hato}, \emph{gayo}
\end{itemize}
\end{tcolorbox}

\subsection*{Your turn}
\begin{itemize}
\raggedright
  \item \emph{Answer:} \textbf{\gu{કાલે} \gu{તમે} \gu{ક્યાં} \gu{હતા}?} — \emph{kāle tame kyā̃ hatā?} — and say which day you took \textbf{\gu{કાલે}} to mean
  \item \emph{Answer:} \textbf{\gu{તમે} \gu{ગઈકાલે} \gu{ઘરે} \gu{હતા} \gu{કે} \gu{નહીં}?} — \emph{tame gaīkāle ghare hatā ke nahī̃?}
\end{itemize}

\begin{tcolorbox}[breakable,colback=teal!4,colframe=teal!35!black,title={Before you move on}]
Which forms in the day agreed with a woman? (\textbf{\gu{હતી}}, \textbf{\gu{ગઈ}}.) Which agreed with the brother? (\textbf{\gu{આવ્યો}}, \textbf{\gu{મોડો}}.)
\end{tcolorbox}
